\chapter{الناسخ والمنسوخ بدلالة السنة والإجماع وأحكام اللعان}

\section{الناسخ والمنسوخ الذي تدل عليه السنة والإجماع}
الحمد لله والصلاة والسلام على رسول الله، أما بعد؛ فقد شرع الإمام الشافعي رحمه الله تعالى في بيان قسم آخر وهو: "الناسخ والمنسوخ الذي تدل عليه السنة والإجماع". ومقصوده أن هناك آيات وأحاديث دلت السنة المطهرة والإجماع على أنها منسوخة.
ويجب التنبيه هنا إلى قاعدة أصولية هامة: أن الإجماع في ذاته لا يَنْسَخ ولا يُنْسَخ؛ لأنه لا إجماع إلا بعد موت رسول الله صلى الله عليه وسلم وانقطاع الوحي. وإنما الإجماع يَدُل على النسخ، أي يُجمع العلماء على أن هذا النص قد نُسخ بنص آخر كان مستقراً في عهد النبوة.

\subsection{نسخ الوصية للوارث بآيات المواريث}
قال الله تبارك وتعالى: ﴿كُتِبَ عَلَيْكُمْ إِذَا حَضَرَ أَحَدَكُمُ الْمَوْتُ إِن تَرَكَ خَيْرًا الْوَصِيَّةُ لِلْوَالِدَيْنِ وَالْأَقْرَبِينَ﴾، وقال في حق الزوجة: ﴿وَالَّذِينَ يُتَوَفَّوْنَ مِنكُمْ وَيَذَرُونَ أَزْوَاجًا وَصِيَّةً لِّأَزْوَاجِهِم﴾. فهاتان الآيتان في سورة البقرة تثبتان الوصية للوالدين وللزوجة.
ثم أنزل الله عز وجل آيات المواريث في سورة النساء مفصلة لأنصبة الورثة. فاحتمل الأمر أن تكون الوصية ثابتة مع الميراث، واحتمل أن تكون آيات المواريث ناسخة لآيات الوصية.
فطلب أهل العلم الدلالة من كتاب الله فلم يجدوا نصاً يقطع بالمراد، فطلبوه في سنة رسول الله صلى الله عليه وسلم، فوجدوا أن أهل المغازي والفتيا نقلوا عامة عن عامة أن النبي صلى الله عليه وسلم قال في خطبته عام الفتح: «لا وصية لوارث». وأجمع أهل العلم قاطبة على ذلك. فدل هذا الحديث مع الإجماع المذكور على أن آيات النساء ناسخة لآيات البقرة، وأن الميراث ناسخ للوصية بالنسبة للوارث.

\subsection{حكم الوصية لغير الوارث ولغير الأقارب}
إذا بطلت الوصية للوارث، فهل تجوز لغير الوارث من الأقارب وغيرهم؟ ذهب الإمام طاوس رحمه الله إلى القول بأن الوصية نسخت للوالدين وثبتت للقرابة غير الوارثين فقط، فمن أوصى لغير قرابة لم تجز وصيته.
فرد الشافعي هذا القول بالدليل، واستدل بحديث عمران بن حصين رضي الله عنه: أن رجلاً من الأنصار أعتق ستة مماليك له عند موته ليس له مال غيرهم، فجزأهم النبي صلى الله عليه وسلم أثلاثاً، فأقرع بينهم، فأعتق اثنين وأرق أربعة.
ووجه الدلالة من الحديث: أن إعتاقه لهم عند الموت وصية، وهم مماليك ليسوا بأقارب له، فدل على جواز الوصية لغير الأقارب، ودل أيضاً على إبطال الاستسعاء في حقهم، وأن الوصية مقتصرة شرعاً على حدود "الثلث" فقط (لأنه أعتق الثلث وأرق الباقي). 

\section{الفرائض التي أنزلها الله نصاً (القذف واللعان)}
شرع الشافعي بعد ذلك في بيان باب الفرائض التي نص عليها القرآن ولا تحتاج إلى بيان إضافي من السنة لكفايتها، ومثّل لذلك بحد القذف. قال تعالى: ﴿وَالَّذِينَ يَرْمُونَ الْمُحْصَنَاتِ ثُمَّ لَمْ يَأْتُوا بِأَرْبَعَةِ شُهَدَاءَ فَاجْلِدُوهُمْ ثَمَانِينَ جَلْدَةً﴾. والمحصنات هنا هن: الحرائر البوالغ العفائف.
فهذا حكم رمي الأجنبي للمرأة بالزنا، أنه يُجلد ثمانين جلدة إن لم يأت بأربعة شهداء، ولا يقبل منه غير ذلك.

\subsection{أحكام اللعان وما زادته السنة على القرآن}
أما إن كان القاذف هو الزوج، فقد فرق القرآن بين حكمه وحكم الأجنبي، فجعل للزوج مخرجاً باللعان: ﴿وَالَّذِينَ يَرْمُونَ أَزْوَاجَهُمْ وَلَمْ يَكُن لَّهُمْ شُهَدَاءُ إِلَّا أَنفُسُهُمْ فَشَهَادَةُ أَحَدِهِمْ أَرْبَعُ شَهَادَاتٍ بِاللَّهِ إِنَّهُ لَمِنَ الصَّادِقِينَ﴾ الآيات.
فإذا تلاعنا، درأ ذلك عنهما الحد. وقد بيّن الشافعي أن القرآن اكتفى ببيان صفة اللعان وعدد الأيمان (أربع شهادات والخامسة اللعنة أو الغضب)، فلم يحتج الصحابة لنقل لفظ النبي صلى الله عليه وسلم في تلقينهم الأيمان لكونها مفصلة وواضحة في التنزيل. 
لكنهم نقلوا في المقابل الأحكام التي زادتها السنة ولم تُذكر في القرآن، وهي: التفريق بين المتلاعنين، ونفي الولد عن الزوج وإلحاقه بالأم. دلالة على أن السنة تستقل بتشريع الأحكام.

\section{الحكم بالظاهر لا بالبواطن والسرائر}
من أعظم القواعد القضائية والمنهجية التي أرساها النبي صلى الله عليه وسلم في قصة المتلاعنين (عويمر العجلاني وزوجته): الحكم بالظاهر. 
فقد قال صلى الله عليه وسلم في شأنهما: «إن أحدكما كاذب»، وذكر علامات خِلْقية ظاهرة في الولد، فقال: «إن جاءت به أُحَيْمِرَ كأنه وحَرَةٌ... فلا أراه إلا قد صدق عليها»، فلما جاءت به على النعت المكروه (الذي يثبت التهمة)، قال صلى الله عليه وسلم: «لولا الأيمان لكان لي ولها شأن»، أي لولا أن الله حكم بدرء الحد عنها باللعان لأقام عليها الحد بسبب وجود الشبه بالرجل المتهَم.

فلم يحكم النبي صلى الله عليه وسلم بالعلامة الظاهرة ولا بما يعلمه من القرائن، بل أمضى حكم الله النافذ باللعان، مغلباً للظاهر الشرعي.
وهذا يوجب على الحكام والقضاة ألا يقضوا إلا بالظاهر (بالبينة أو الإقرار)، ولا يقضوا بعلمهم الشخصي ولا بالقرائن المجردة عن البينة. كما قال صلى الله عليه وسلم في الحديث الصحيح: «إنما أنا بشر، وإنكم تختصمون إلي، ولعل بعضكم أن يكون ألحن بحجته من بعض، فأقضي له على نحو ما أسمع، فمن قضيت له من حق أخيه شيئاً فلا يأخذه، فإنما أقطع له قطعة من النار». 
فالشريعة جاءت لتبني الأحكام على الظواهر المنضبطة، وتَكِل السرائر إلى الخبير بها سبحانه وتعالى.